% 連続体力学 連続体力学-1 問題4（004-accel） 解いた日 2026-10-17
%% kotae
\toi{1}{△}{記述}
\[ \vect{a}=\frac{\mathrm{D}\vect{u}}{\mathrm{D}t}=\pdv{\vect{u}}{t}+(\vect{u}\cdot\vect{\nabla})\vect{u} \]
\nuke{式(4)で $A=u_{j}$}
\[ a_{j}=\pdv{u_{j}}{t}+\sum_{i=1}^{3}u_{i}\pdv{u_{j}}{x_{i}} \]
\naze{示すべき式をそのまま書いただけで導出になっていない。式(2)の第 $j$ 成分 $a_{j}=\mathrm{D}u_{j}/\mathrm{D}t$ をとり，$u_{j}$ を1つのスカラー量とみなしてスカラーについての式(4)を $A=u_{j}$ として当てはめる，という段が要の論点である。}
\toi{2}{○}{記述}
\begin{align*}
\vect{a}&=0+(\vect{u}\cdot\vect{\nabla})\vect{u}\\
&=\Bigl(-\Omega x_{2}\pdv{}{x_{1}}+\Omega x_{1}\pdv{}{x_{2}}\Bigr)\vect{u}\\
&=-\Omega x_{2}(0,\Omega,0)+\Omega x_{1}(-\Omega,0,0)\\
&=(-\Omega^{2}x_{1},-\Omega^{2}x_{2},0)
\end{align*}
（図：$x_{1}x_{2}$ 平面に，点から原点に向かう矢印）
\begin{align*}
\abs{\vect{a}}&=\Bigl|-\Omega^{2}\sqrt{x_{1}^{2}+x_{2}^{2}}\Bigr|\\
&=\ayamari{-r\Omega^{2}}{r\Omega^{2}}
\end{align*}
\nuke{向き：$x_{3}$ 軸に向かう向き}
\naze{向きは図の矢印だけで，文で答えていない。$\vect{a}=-\Omega^{2}(x_{1},x_{2},0)$ は $x_{3}$ 軸に垂直に軸へ向かう向き，と書く。大きさは絶対値なので正の $\Omega^{2}r$ である。}
\toi{3}{△}{立式}
\[ -\Omega x_{2}=\frac{\mathrm{D}x_{1}}{\mathrm{D}t},\quad \Omega x_{1}=\frac{\mathrm{D}x_{2}}{\mathrm{D}t} \]
\[ x_{2}=-\frac{1}{\Omega}\frac{\mathrm{D}x_{1}}{\mathrm{D}t},\quad \Omega x_{1}=\frac{\mathrm{D}}{\mathrm{D}t}\Bigl(-\frac{1}{\Omega}\frac{\mathrm{D}x_{1}}{\mathrm{D}t}\Bigr) \]
\[ \therefore\ \frac{\mathrm{D}^{2}x_{1}}{\mathrm{D}t^{2}}=-\Omega^{2}x_{1} \]
\[ x_{1}(t)=C_{1}\sin(\Omega t)+C_{2}\cos(\Omega t) \]
$x_{1}(0)=X_{1}$ より $C_{2}=X_{1}$
\begin{align*}
x_{2}(t)&=-\frac{1}{\Omega}\bigl(C_{1}\Omega\cos(\Omega t)-C_{2}\Omega\sin(\Omega t)\bigr)\\
&=-C_{1}\cos(\Omega t)+C_{2}\sin(\Omega t)
\end{align*}
$x_{2}(0)=X_{2}$ より $C_{1}=-X_{2}$
\[ \vect{x}(t)=\begin{pmatrix}-X_{2}\sin(\Omega t)+X_{1}\cos(\Omega t)\\X_{2}\cos(\Omega t)+X_{1}\sin(\Omega t)\\\ayamari{0}{X_{3}}\end{pmatrix} \]
\naze{第3成分は $\mathrm{D}x_{3}/\mathrm{D}t=u_{3}=0$ と初期条件 $x_{3}(0)=X_{3}$ から $x_{3}=X_{3}$ である。0 では時刻0に $\vect{X}$ にいたことに反する。}
\begin{align*}
\frac{\mathrm{D}^{2}x_{1}}{\mathrm{D}t^{2}}&=\Omega^{2}X_{2}\sin(\Omega t)-\Omega^{2}X_{1}\cos(\Omega t)\\
&=-\Omega^{2}\bigl(-X_{2}\sin(\Omega t)+X_{1}\cos(\Omega t)\bigr)\\
&=-\Omega^{2}x_{1}=a_{1}
\end{align*}
\begin{align*}
\frac{\mathrm{D}^{2}x_{2}}{\mathrm{D}t^{2}}&=-\Omega^{2}X_{2}\cos(\Omega t)-\Omega^{2}X_{1}\sin(\Omega t)\\
&=-\Omega^{2}\bigl(X_{2}\cos(\Omega t)+X_{1}\sin(\Omega t)\bigr)\\
&=-\Omega^{2}x_{2}=a_{2}
\end{align*}
\toi{4}{○}{記述}
\[ \vect{a}=\frac{\mathrm{D}\vect{u}}{\mathrm{D}t}=\pdv{\vect{u}}{t}+u_{1}\pdv{\vect{u}}{x_{1}} \]
\[ a_{1}=-\frac{x_{1}}{(t+t_{0})^{2}}+\frac{x_{1}}{t+t_{0}}\cdot\frac{1}{t+t_{0}}=0 \]
\[ a_{2}=a_{3}=0 \]
\[ \vect{u}=\frac{\mathrm{D}\vect{x}}{\mathrm{D}t},\quad u_{1}=\frac{x_{1}}{t+t_{0}}=\frac{\mathrm{D}x_{1}}{\mathrm{D}t} \]
\[ \therefore\ \frac{\mathrm{D}t}{t+t_{0}}=\frac{\mathrm{D}x_{1}}{x_{1}},\quad \therefore\ \ln(t+t_{0})=\ln x_{1}+C \]
\[ \therefore\ x_{1}=C(t+t_{0}) \]
$t=0$ における $x$ を $X_{1}$ とすると $C=X_{1}/t_{0}$
\[ x_{1}=\frac{X_{1}}{t_{0}}(t+t_{0}) \]
（図：等間隔に並んだ点の列が，右へ行くほど間隔の広がった点の列に変わる。上に左向きと右向きの矢印）

固定しているところから見ると流れは遅くなっているように見えるが，粒子は等速直線運動をしているから。\nuke{遅い粒子が入れ替わる}
\naze{各粒子が等速であることに加えて，固定点の速度 $x_{1}/(t+t_{0})$ が減るのは，速度 $X_{1}/t_{0}$ の小さい（出発点が原点に近い）粒子ほど遅れてその点に着き，次々と入れ替わるからである，まで書くと理由になる。}
%% kizuki
\item 式(4)はスカラー量 $A$ についての関係である。ベクトルの加速度に使うときは「第 $j$ 成分 $a_{j}=\mathrm{D}u_{j}/\mathrm{D}t$ で $A=u_{j}$ とおく」と1行書けば導出になる（(1)）。
\item 軌道を求めたら，最後に $t=0$ を代入して $\vect{x}(0)=\vect{X}$ になるかを3成分とも確かめる。$u_{3}=0$ の成分は「動かない」なので $x_{3}=X_{3}$ で，0 ではない（(3)）。大きさは絶対値なので負にならないことも見直す（(2)）。
\item (4) の理由は問題1 (4) と同じ形である。粒子ごとに見た量（ここでは速度 $X_{1}/t_{0}$）が一定でも，固定点には別の粒子が次々と来るので場の値は変わる。「どの粒子が来るか」を書く。
